% Build from the repository root:
% xelatex -output-directory=reports reports/phase2.tex
% Requires XB Niloofar and the Latin Modern fonts.
\documentclass[12pt,a4paper]{article}
\usepackage[a4paper,top=23mm,bottom=22mm,left=26mm,right=26mm]{geometry}
\usepackage{amsmath,booktabs,graphicx,pgf,fontspec,unicode-math}
\usepackage[bidi=default]{babel}
\babelprovide[import,main]{persian}
\babelprovide[import]{english}
\babelfont[persian]{rm}{XB Niloofar}
\babelfont[english]{rm}{Latin Modern Roman}
\setmathfont{Latin Modern Math}

\pagestyle{empty}
\setlength{\parindent}{0pt}
\setlength{\parskip}{4pt}
\setlength{\emergencystretch}{1em}
\linespread{1.16}
\renewcommand{\arraystretch}{1.2}

\newcommand{\en}[1]{\foreignlanguage{english}{#1}}
\newcommand{\m}[1]{\babelsublr{\mbox{\ensuremath{#1}}}}
\newcommand{\fa}[1]{\leavevmode\hbox{\beginR #1\endR}}
\newcommand{\topic}[2]{%
  \par\addvspace{12pt}%
  {\large\bfseries #1\enspace #2\par}%
  \nopagebreak[4]\vspace{1pt}%
}
\newcommand{\nacc}[1]{\hfill{\small\mdseries ج\m{#1}}}
\newcommand{\legend}[3]{%
  \par\nopagebreak[4]%
  {\small\textbf{#1 \m{#2}.}\enspace #3\par}%
}
\providecommand{\mathdefault}[1]{#1}
\newcommand{\plot}[2]{%
  \par\nopagebreak[4]\vspace{3pt}%
  \noindent\makebox[\linewidth][c]{%
    \en{\resizebox{#1\linewidth}{!}{\input{figures/phase2/latex/#2.pgf}}}%
  }\par%
}

\begin{document}

{\small گزارش فاز دوم پروژه آمار\par}
\vspace{5pt}
{\fontsize{25}{36}\selectfont\bfseries احتمال و توزیع‌های احتمالی\par}
\smallskip
\hrule height 0.5pt
\medskip
\textbf{سوال پروژه}\enspace
میانگین ضخامت قشر انتورینال در دو نیمکره چه رابطه‌ای با مجموع نمره حوزه‌های
\en{CDR} دارد و الگوی توزیع این دو متغیر در گروه‌های مختلف وضعیت شناختی
چه تفاوت‌هایی نشان می‌دهد؟

\topic{\m{1.}}{تعریف رویدادها و متغیرهای تصادفی\nacc{5}}
در هر یک از فایل‌های \en{UDS} و \en{MRI}، برای \m{1237} نفر یک رکورد
به ازای هر شناسه \en{NACCID} وجود داشت. بعد از اتصال یک‌به‌یک، نمره،
وضعیت شناختی و ضخامت هر دو نیمکره برای همه افراد معتبر بود و کسی حذف نشد.
برای هر فرد همان رکورد موجود در هر فایل را به کار بردیم. شماره و تاریخ مراجعه
در دسترس نبود، پس فاصله زمانی اسکن و ارزیابی بالینی معلوم نیست.
فایل \en{UDS} پس از جای‌گذاری داده‌های گمشده به دست ما رسیده و معلوم نیست
کدام مقدارها جای‌گذاری شده‌اند.

یک نفر را با احتمال برابر از نمونه انتخاب می‌کنیم. \m{X} میانگین ضخامت قشر
انتورینال در دو نیمکره بر حسب میلی‌متر و \m{Y} نمره \en{CDRSUM} همان فرد
است. این نمره از جمع نمره شش بخش ارزیابی شناخت و عملکرد به دست می‌آید.

رویداد \m{A=\{Y>0\}} یعنی نمره بیشتر از صفر و رویداد \m{B=\{X\leq 3\}}
یعنی ضخامت حداکثر \m{3} میلی‌متر. این مرز از گرد کردن میانگین ضخامت در
فاز اول به دست آمده و کاربردش در این تحلیل فقط گروه‌بندی توصیفی است.

\topic{\m{2.}}{احتمال‌های حاشیه‌ای و مشترک\nacc{3}}
\legend{جدول}{1}{تعداد افراد بر اساس ضخامت و نمره}
\begin{center}
\begin{tabular*}{\linewidth}{@{\extracolsep{\fill}}ccc r@{}}
\toprule
\fa{کل} & \fa{نمره صفر} & \fa{نمره بیشتر از صفر} & \fa{گروه ضخامت}\\
\midrule
\m{537}  & \m{254} & \m{283} & \fa{حداکثر \m{3} میلی‌متر}\\
\m{700}  & \m{503} & \m{197} & \fa{بیشتر از \m{3} میلی‌متر}\\
\midrule
\m{1237} & \m{757} & \m{480} & \fa{کل}\\
\bottomrule
\end{tabular*}
\end{center}

در کل نمونه، \m{P(A)\approx38.80\%} و \m{P(B)\approx43.41\%} است.
از \m{1237} نفر، \m{283} نفر هر دو ویژگی را دارند. در نتیجه،
\m{P(A\cap B)\approx22.88\%} است.

\topic{\m{3.}}{احتمال شرطی\nacc{3}}
در گروه ضخامت حداکثر \m{3} میلی‌متر، \m{283} نفر از \m{537} نفر نمره
بالاتر از صفر دارند، یعنی \m{52.70\%}. این سهم در گروه ضخامت بیشتر،
\m{197} نفر از \m{700} نفر، یعنی \m{28.14\%} است. اختلاف دو گروه
\m{24.56} واحد درصد است.

\topic{\m{4.}}{بررسی مفهوم استقلال}
اگر دو رویداد مستقل بودند، احتمال مشترک با حاصل‌ضرب احتمال‌های آن‌ها
برابر بود. احتمال مشترک حدود \m{22.88\%} و حاصل‌ضرب حدود \m{16.85\%} است.
پس دو رویداد در این نمونه وابسته‌اند.

\clearpage
\topic{\m{5.}}{توزیع گسسته و مدل دوجمله‌ای}
فرض می‌کنیم \m{20} بار از همین نمونه یک نفر را با احتمال برابر انتخاب
می‌کنیم و بعد او را به مجموعه برمی‌گردانیم. انتخاب‌ها مستقل‌اند و احتمال
نمره بالاتر از صفر در هر بار \m{p=480/1237} است. تعداد انتخاب‌ها
با نمره بالاتر از صفر، \m{K}، توزیع دوجمله‌ای دارد و
\[
P(K=8)=\binom{20}{8}p^8(1-p)^{12}\approx17.86\%.
\]
بدون جای‌گذاری، احتمال هر بار با انتخاب‌های قبلی تغییر می‌کند.
چون \m{20} نفر بخش کوچکی از نمونه \m{1237} نفری‌اند، مدل دوجمله‌ای
برای آن حالت هم تقریب معقولی است.

\topic{\m{6.}}{متغیر پیوسته و توزیع نرمال\nacc{2}}
ضخامت بیشتر افراد نزدیک \m{3} میلی‌متر است و دنباله توزیع به سمت
ضخامت‌های کمتر کشیده شده است. منحنی نرمال را با میانگین و واریانس
همین اندازه‌ها رسم کردیم.

\plot{0.99}{thickness_normal}
\legend{شکل}{1}{سمت چپ، هیستوگرام ضخامت و منحنی نرمال با چگالی بر حسب
معکوس میلی‌متر. سمت راست، چندک‌های مشاهده‌شده در برابر چندک‌های نرمال.}

چندک‌های میانی به خط برابری نزدیک‌اند، اما در دو سر نمودار از آن فاصله
می‌گیرند. تقریب نرمال برای بخش میانی توزیع مناسب‌تر است.

\topic{\m{7.}}{امید ریاضی و واریانس\nacc{2}}
برای انتخاب یک نفر از این نمونه، امید ریاضی همان میانگین جدول است.
واریانس تجربی را با مخرج \m{n} محاسبه کردیم. مخرج انحراف معیار نمونه
در فاز اول \m{n-1} بود.

واریانس ضخامت بر حسب میلی‌متر مربع است و ریشه آن، حدود \m{0.36} میلی‌متر،
مقیاس پراکندگی ضخامت‌ها دور میانگین را نشان می‌دهد. واریانس نمره
بر حسب مربع واحد نمره است.

\legend{جدول}{2}{خلاصه توزیع نمره و ضخامت در کل نمونه}
\begin{center}
\begin{tabular*}{\linewidth}{@{\extracolsep{\fill}}cccc r@{}}
\toprule
\fa{دامنه} & \fa{میانه} & \fa{واریانس} & \fa{میانگین} & \fa{متغیر}\\
\midrule
\fa{\m{0} تا \m{15}}             & \m{0}     & \m{2.966} & \m{0.822} & \fa{نمره \en{CDRSUM}}\\
\fa{\m{1.214} تا \m{3.925}}      & \m{3.058} & \m{0.127} & \m{3.014} & \fa{ضخامت به میلی‌متر}\\
\bottomrule
\end{tabular*}
\end{center}

\clearpage
نمره \m{61.2\%} افراد صفر است. نمره‌های دنباله راست میانگین را به
\m{0.822} واحد رسانده‌اند و بیشتر پراکندگی نمره از همین مقادیر می‌آید.
با قاعده \m{1.5} برابر فاصله بین چارک‌ها، \m{115} نمره بالا و
\m{22} ضخامت پایین علامت خوردند. همه این مقادیر در تحلیل باقی ماندند.

\plot{0.85}{cdrsum}
\legend{شکل}{2}{تعداد افراد در هر نمره مشاهده‌شده.}

\topic{\m{8.}}{دو متغیر تصادفی و کوواریانس\nacc{4}}
برای هر فرد، فاصله نمره و ضخامت را از میانگین‌هایشان در هم ضرب کردیم.
میانگین این ضرب‌ها، یعنی کوواریانس، حدود \m{-0.179} شد.
علامت منفی یعنی در مجموع، ضخامت کمتر با نمره بالاتر همراه است.
در شکل، حتی افرادی با ضخامت نزدیک به هم نمره‌های متفاوتی دارند.

\plot{0.85}{covariance}
\legend{شکل}{3}{رابطه نمره و ضخامت. هر نقطه یک نفر را نشان می‌دهد.}

\topic{\m{9.}}{قضیه بیز}
در کل نمونه، احتمال نمره بالاتر از صفر \m{38.80\%} است.
اگر بدانیم ضخامت فرد حداکثر \m{3} میلی‌متر است، با قضیه بیز به همان
احتمال شرطی بخش \m{3} می‌رسیم.
\[
P(A\mid B)=\frac{P(B\mid A)P(A)}{P(B)}
          =\frac{(283/480)(480/1237)}{537/1237}\approx52.70\%.
\]

\clearpage
\topic{ج.\m{1}}{احتمال شرطی وضعیت شناختی}
وضعیت شناختی را از \en{NACCUDSD} در فایل بالینی گرفتیم. درصدها در
هر ستون نسبت به تعداد افراد همان گروه ضخامت محاسبه شده‌اند.
مخرج‌ها به ترتیب \m{537} و \m{700} هستند.

\legend{جدول}{3}{وضعیت شناختی در دو گروه ضخامت و میانه نمره و ضخامت در هر وضعیت}
\begin{center}
\begin{tabular*}{\linewidth}{@{\extracolsep{\fill}}cccc r@{}}
\toprule
\shortstack{\fa{میانه ضخامت}\\[2pt]\fa{به میلی‌متر}} &
\fa{میانه نمره} &
\shortstack{\fa{ضخامت بیشتر از}\\[2pt]\fa{\m{3} میلی‌متر}} &
\shortstack{\fa{ضخامت حداکثر}\\[2pt]\fa{\m{3} میلی‌متر}} &
\fa{وضعیت شناختی}\\
\midrule
\m{3.103} & \m{0}   & \m{555\;(79.29\%)} & \m{292\;(54.38\%)} & \fa{شناخت طبیعی}\\
\m{3.188} & \m{1}   & \m{6\;(0.86\%)}    & \m{1\;(0.19\%)}    & \fa{اختلال بدون \en{MCI}}\\
\m{2.966} & \m{1}   & \m{111\;(15.86\%)} & \m{141\;(26.26\%)} & \fa{اختلال شناختی خفیف}\\
\m{2.691} & \m{4.5} & \m{28\;(4.00\%)}   & \m{103\;(19.18\%)} & \fa{دمانس}\\
\bottomrule
\end{tabular*}
\end{center}

برای نمونه، دمانس در \m{19.18\%} افراد گروه ضخامت حداکثر \m{3} میلی‌متر
و \m{4.00\%} افراد گروه ضخامت بیشتر ثبت شده است.

\plot{0.99}{by_status}
\legend{شکل}{4}{توزیع ضخامت در سمت چپ و نمره در سمت راست برای چهار وضعیت شناختی.
خط داخل هر جعبه میانه است. کدهای \m{1} تا \m{4} به ترتیب سطرهای جدول
\m{3} هستند.}

\topic{}{جمع‌بندی فاز دوم}
در این نمونه، نمره بالاتر از صفر در گروه با ضخامت کمتر رایج‌تر بود.
میانه ضخامت در گروه دمانس
\m{2.69} میلی‌متر و در گروه شناخت طبیعی \m{3.10} میلی‌متر است.
میانه نمره در این دو گروه به ترتیب \m{4.5} و صفر است. ضخامت افراد
گروه‌های مختلف همپوشانی دارد و در گروه اختلال بدون \en{MCI} فقط
\m{7} نفر هستند.

نمونه‌گیری تصادفی نبوده و داده‌ها مشاهده‌ای‌اند. تفسیر این الگوها
به رابطه آماری در همین نمونه محدود است.

\end{document}
